% !TEX root = ../main.tex
\section{SQL 틀}
\label{ch:appendix-sql}

이 부록은 두 분석의 그래프가 바탕으로 삼는 질의 틀 세 가지를 줄여서 보여 줘요. 전체 파일은 실험 폴더의 \texttt{sql/}에 있어요(부록~\ref{sec:pipeline-stages}). \texttt{\_\_DEST\_\_}처럼 밑줄 두 개로 감싼 자리 표시는 질의를 돌리기 전에 \texttt{scripts/run\_bq.py}가 채워 넣어요. 그리고 \texttt{@}로 시작하는 이름은 질의 매개변수(질의에 밖에서 넣어 주는 값)예요.

\dissertationSubheading[sec:sql-materialize]{로그를 미리 표로 만들어 두는 틀}

스캔~A와 스캔~B는 같은 질의(\texttt{sql/01\_materialize\_logs.sql})를 서로 다른 기간에 대해 돌려요. 이 질의는 공개 로그 표를 한 번만 읽어요. 그리고 로그마다 종류 꼬리표를 붙이고, 이 연구에 필요한 네 종류를 달마다 나눈 프로젝트 표에 남겨요. 토픽(topic, 로그에 붙는 꼬리표) 해시와 GMX 갈래, Aave 갈래를 줄여서 보이면 다음과 같아요:

\begin{verbatim}[baselinestretch=1]
CREATE TABLE `__DEST__`
PARTITION BY TIMESTAMP_TRUNC(block_timestamp, MONTH) CLUSTER BY kind AS
WITH tagged AS (
  SELECT *, CASE
      WHEN ARRAY_LENGTH(topics) = 3 AND topics[OFFSET(0)] = '0x8c5b...b925'
        THEN 'A'                                  -- ERC-20 Approval
      WHEN ARRAY_LENGTH(topics) = 3 AND topics[OFFSET(0)] = '0xddf2...b3ef'
        THEN 'T'                                  -- ERC-20 Transfer
      WHEN ... THEN 'G'                           -- GMX V2 PositionDecrease
      WHEN ... THEN 'V'                           -- Aave V3 pool events
    END AS kind
  FROM `bigquery-public-data.goog_blockchain_arbitrum_one_us.logs`
  WHERE block_timestamp >= TIMESTAMP(@start_ts)
    AND block_timestamp < TIMESTAMP(@end_ts))
SELECT block_timestamp, block_number, log_index, kind,
  FROM_HEX(SUBSTR(address, 3)) AS emitter,
  IF(kind IN ('A', 'T'), FROM_HEX(SUBSTR(topics[OFFSET(1)], 27)), NULL) AS a,
  IF(kind IN ('A', 'T'), FROM_HEX(SUBSTR(topics[OFFSET(2)], 27)), NULL) AS b,
  IF(kind IN ('A', 'T'), SAFE.FROM_HEX(SUBSTR(data, 3)), NULL) AS val -- ...
FROM tagged WHERE kind IS NOT NULL
\end{verbatim}

허락에서는 \texttt{a}가 주인이고 \texttt{b}가 스펜더예요. 송금에서는 \texttt{a}가 보낸 쪽이고 \texttt{b}가 받는 쪽이에요. 두 경우 모두 \texttt{emitter}는 토큰 계약이고, \texttt{val}은 32바이트 금액이에요. 그리고 주소는 20바이트 값으로 저장해요.

\dissertationSubheading[sec:sql-ego]{에고 네트워크 뽑아내기 틀}

\texttt{sql/03\_ego\_events.sql} 질의는 한 기간에 대해 미리 만들어 둔 로그 표를 읽어요. 그리고 한쪽 끝에 계획의 코호트에 든 계약이 있는 허락과 송금만 남겨요. 계획의 코호트는 \texttt{contract\_rep.cohort} 표에 들어 있어요. 허락 부분을 줄여서 보이면 다음과 같아요:

\begin{verbatim}[baselinestretch=1]
CREATE OR REPLACE TABLE `dissertation-bq.contract_rep.approvals___SUFFIX__`
PARTITION BY TIMESTAMP_TRUNC(block_timestamp, MONTH) AS
SELECT l.block_timestamp, l.block_number, l.log_index,
       l.emitter AS token_address, l.a AS owner, l.b AS spender,
       u256(l.val) AS value
FROM `dissertation-bq.contract_rep.__SRC__` AS l
WHERE l.kind = 'A'
  AND l.block_timestamp >= TIMESTAMP(@start_ts)
  AND l.block_timestamp < TIMESTAMP(@end_ts)
  AND (l.a IN (SELECT address FROM `dissertation-bq.contract_rep.cohort`)
       OR l.b IN (SELECT address FROM `dissertation-bq.contract_rep.cohort`));
-- transfers___SUFFIX__: the same with kind = 'T', from_address = a,
-- to_address = b and zero amounts dropped: AND LTRIM(TO_HEX(l.val), '0') != ''
\end{verbatim}

파일 맨 위에 정의한 함수 \texttt{u256}은 32바이트 금액을 원래 단위(소수점 조정을 하지 않은 가장 작은 단위)의 배정밀도 실수(double)로 읽어요. 금액이 0인 허락은 남겨요. 금액이 0인 허락은 허락을 0으로 만드는 일이기 때문이에요. 반면 금액이 0인 송금은 버려요. 수정 분석(등록 뒤에 금액을 달러로 바꿔 다시 한 분석)은 이 질의를 다시 돌리지 않아요. 그 대신 \texttt{sql/12\_usd\_ego.sql}이 이 표들에서 다섯 토큰의 이벤트를 남겨요. 남기는 것은 스캔~A에서 다섯 토큰에 대해 주인 세 명 이상이 허락한 계약이 한쪽 끝에 있는 이벤트예요. 그리고 송금마다 그 블록이 있는 UTC 날짜의 토큰 가격으로 금액을 매겨요. 그다음 \texttt{sql/21\_usd12m\_cohorts.sql}이 각 동결일(점수를 고정하는 날)의 코호트를 그 동결일의 점수 기간(점수를 계산하는 데 쓰는 기간)에 있는 허락에서 뽑아요.

\dissertationSubheading[sec:sql-pairs]{마지막 허락과 쌍 표를 만드는 틀}

수정 분석에서는 \texttt{sql/14\_usd\_anchor\_pairs.sql} 질의가 한 기준 시점(앵커)까지의 에고 이벤트를 두 그래프의 화살표 목록으로 바꿔요. 등록 분석(미리 등록한 그대로 한 분석)에서는 \texttt{sql/04\_anchor\_pairs.sql} 질의가 같은 일을 원래 단위로, $b=1$을 쓰고 상한 없이 했어요. 줄여서 보이면 다음과 같아요:

\begin{verbatim}[baselinestretch=1]
CREATE OR REPLACE TABLE `dissertation-bq.contract_rep.u_allow_triples___TAG__` AS
SELECT r.token_address, r.owner, r.spender, r.amount * p.price_usd AS usd,
       TIMESTAMP_DIFF(TIMESTAMP(@anchor_ts), r.block_timestamp, MICROSECOND)
         / 8.64e10 AS age_days
FROM (
  SELECT *, ROW_NUMBER() OVER (
           PARTITION BY token_address, owner, spender
           ORDER BY block_number DESC, log_index DESC) AS rn
  FROM `dissertation-bq.contract_rep.u_approvals`
  WHERE block_timestamp >= TIMESTAMP(@start_ts)
    AND block_timestamp <= TIMESTAMP(@anchor_ts)
) AS r
JOIN `dissertation-bq.contract_rep.u_prices` AS p
  ON p.token = r.token_address AND p.day = DATE(TIMESTAMP(@anchor_ts))
WHERE r.rn = 1 AND r.value > 0;

CREATE OR REPLACE TABLE `dissertation-bq.contract_rep.u_allow_pairs___TAG__` AS
SELECT owner, spender,
       SUM(sig(age_days) * vt(usd, @b_allowance)) AS w_er,  -- EndorseRank
       SUM(LEAST(usd, @cap_allowance)) AS w_raw,    -- C-PR allowance layer
       MAX(sig(age_days)) AS max_sig
FROM `dissertation-bq.contract_rep.u_allow_triples___TAG__`
GROUP BY owner, spender;
-- u_transfer_pairs___TAG__: per (from_address, to_address) up to the anchor,
-- w_awp = SUM(sig * vt(usd, @b_transfer)), w_raw = SUM(usd * sig), MAX(sig)
\end{verbatim}

파일 맨 위에 정의한 함수 \texttt{sig}와 \texttt{vt}는 식~\eqref{eq:logistic-decay}의 로지스틱 감쇠(오래된 일일수록 가볍게 세는 무게)와 식~\eqref{eq:value-transform}의 값 변환(금액이 아무리 커도 1을 넘지 않게 바꾸는 식)이에요. 로지스틱 감쇠에서는 $k=0.01$이고 $t_0=180$일이에요. 값 변환은 $bz=40$부터 값을 1로 정해 두어요. 그 자리부터는 배정밀도로 계산해도 값이 1과 같기 때문이에요. 창 함수(window function, 묶음마다 행에 차례를 매기는 SQL 기능)는 토큰, 주인, 스펜더마다 블록 번호와 로그 번호(log index)로 마지막 허락을 골라요. 이 마지막 허락은 금액이 0보다 클 때만 남겨요. 그래서 허락 취소(허락을 거두어들임)가 있으면 그 토큰은 그 쌍에서 빠져요. 그리고 마지막 허락에는 기준 시점 날의 가격으로 금액을 매겨요. 1달러당 기울기 $b_A=\ln 3/100$과 $b_T=\ln 3/170$, 그리고 상한 $U=19{,}000$달러는 \texttt{constants.json}에서 매개변수로 넘겨받아요. 쌍 표의 각 행에는 모든 방법의 무게가 함께 들어 있어요. 그래서 한 기준 시점의 모든 점수는 같은 행들에서 나와요. 같은 질의는 기울기를 10배로 키운 경우와 10분의 1로 줄인 경우의 엔도스랭크(EndorseRank) 무게와 AWP 무게도 써 둬요. 표~\ref{tab:robustness-slope}의 금액 기울기 점검은 이것으로 계산해요.
